\chapter{研究方案}

\section{研究方法及实验手段}

本项目采用理论分析、算法设计、模块化实现与闭环仿真验证相结合的研究方法，围绕三个研究内容展开：第一，研究基于预测不确定性的轻量化风险感知运动基元规划；第二，研究通信状态与邻机相关性联合驱动的协同约束自适应方法；第三，构建面向动态复杂环境的无人机规划闭环仿真验证系统，并在统一实验条件下完成前两项方法的对比与消融验证。

需要说明的是，第三个研究内容在论文结构中作为系统构建与综合验证部分放在前两个方法研究之后，但平台基础链路将在研究全过程中同步建设。前期首先打通Gazebo、PX4和ROS2之间的基本感知—规划—控制闭环，为第一项方法提供单机动态环境实验条件；随后扩展多无人机命名空间、状态/轨迹共享和通信扰动接口，为第二项方法提供验证条件；最后将动态障碍、通信扰动、算法切换、指标记录和批量实验能力统一到同一套闭环仿真系统中，形成第三项研究内容。

\subsection{基于预测不确定性的轻量化风险感知运动基元规划方法}

首先利用LiDAR等局部传感器观测获取障碍信息，对连续观测中的动态障碍进行目标提取、数据关联和状态跟踪，得到位置、速度、尺寸及时间戳等状态量。考虑高频局部规划对实时性的要求，短时预测优先采用匀速或局部线性等轻量运动模型，并通过状态协方差传播、预测残差统计或安全包络膨胀等方式描述预测误差随时间的扩展，使预测模块输出未来状态均值$\boldsymbol{\mu}_{o,\tau}$及协方差$\boldsymbol{\Sigma}_{o,\tau}$或等价风险边界。

在规划层采用运动基元滚动评估框架，根据无人机当前速度、加速度、飞行方向和动力学约束生成有限候选轨迹集合。对每条候选轨迹计算目标趋近、平滑性、运动代价和动态风险，其中风险项综合障碍未来状态均值、不确定性范围及候选轨迹之间的时空关系。预测不确定程度增大时，通过扩大风险边界或提高风险代价，使候选轨迹排序对预测可信程度产生响应。

为控制在线计算开销，风险计算优先采用距离、时间、协方差投影或安全包络查询等低复杂度表达，不直接在每个规划周期执行高维概率积分或复杂连续概率优化。在软风险评价之外，保留TTC和最小安全距离等硬安全过滤；当正常候选均不满足安全条件时，触发减速、制动、后退、侧移或受控爬升等应急动作。通过确定性预测基线、不确定性感知方法和相应消融实验，分析新增不确定性处理所带来的安全收益与计算开销。

\subsection{通信状态与邻机相关性联合驱动的协同约束自适应方法}

在单机动态规划基础上，将研究场景扩展至多无人机系统。多机通过共享状态和计划轨迹完成冲突判断，并在消息链路中引入固定时延、随机时延、不同丢包率、短时通信中断及动态链路质量等非理想通信条件。每条邻机状态或轨迹消息记录发送时间、接收时间和最近一次有效更新时间，并由此提取时延、丢包统计、信息年龄和链路状态等通信特征。

通信特征用于描述共享信息的可靠程度，但不直接等价为邻机对当前规划的重要程度。进一步结合机间距离、相对速度、潜在TTC及轨迹交叉趋势等因素评价邻机规划相关性，从而区分“信息存在一定时延但与本机冲突风险较高的关键邻机”和“通信良好但与当前规划关系较弱的普通邻机”。

根据通信可靠程度和邻机相关性构建重点协同邻机集合，并对进入集合的邻机调整轨迹可信时间窗、状态外推范围、机间安全裕度和冲突约束强度。通信状态良好且邻机信息新鲜时充分利用共享轨迹；信息逐渐过期时缩短可信时间窗并提高安全裕度；关键信息严重失效时降低对长期共享轨迹的依赖，以本机感知、短时外推和硬安全约束完成保守避碰。正常协同、保守协同和局部自主作为通信退化条件下的处理状态保留，但不将模式切换本身作为创新核心。

\subsection{面向动态复杂环境的无人机规划闭环仿真验证系统构建}

第三项研究内容面向前两项方法的统一实现、闭环验证和实验复现。系统以Gazebo负责三维静态/动态障碍环境、多无人机场景和LiDAR等传感器仿真，以PX4负责无人机动力学与底层飞行控制，以ROS2承载感知、状态估计、预测、规划、指令仲裁、多机信息共享和通信扰动模拟等上层模块。

系统按照“环境与传感器观测—ROS2感知/预测—局部规划—指令仲裁与接口转换—PX4控制执行—飞行状态反馈”的链路组织单机闭环。多机场景在此基础上采用独立命名空间管理各无人机节点，通过统一的状态和轨迹消息接口进行邻机信息共享，并在通信中间层注入时延、丢包、短时中断等扰动，使通信实验条件可配置、可重复。

针对动态环境实验，构建可参数化的静态和动态障碍场景，支持障碍数量、密度、尺寸、运动速度、转向/变速模式及预测扰动程度调整；针对多机实验，支持无人机数量、初始位置、任务目标和通信参数调节。系统同时记录轨迹、碰撞、最小安全距离、规划时间、控制状态、机间距离和通信统计等数据，并提供统一的算法切换和指标统计接口，以保证基线、改进方法和消融实验在相同条件下进行比较。

该系统研究点的价值主要体现在闭环集成和统一验证，而不是将Gazebo、PX4或ROS2本身作为创新。通过将感知、规划、控制、通信扰动和评价环节统一起来，可避免仅在离线轨迹或理想通信条件下评价算法，并为后续实体无人机迁移保留统一接口。

\section{技术路线及关键技术}

本项目总体技术路线按照“方法研究—多机扩展—系统集成验证”的逻辑展开。首先在单机动态环境中建立短时预测、不确定性表征和风险感知运动基元规划链；随后扩展至多无人机，在状态/轨迹共享基础上加入通信状态评价、邻机相关性判断和协同约束调整；最后在Gazebo、PX4和ROS2环境中将两类方法与感知、控制、通信扰动和评价模块统一集成，形成完整闭环实验系统。

技术流程可概括为：

\begin{center}
动态复杂环境任务需求\\
$\downarrow$\\
LiDAR等局部环境感知与动态障碍状态估计\\
$\downarrow$\\
短时预测与不确定性表征\\
$\downarrow$\\
低复杂度风险计算与运动基元筛选\\
$\downarrow$\\
TTC/安全距离硬安全过滤\\
$\downarrow$\\
单机动态环境规划\\
$\downarrow$\\
多无人机状态/轨迹共享\\
$\downarrow$\\
通信状态评价与邻机相关性判断\\
$\downarrow$\\
重点邻机筛选与协同约束自适应\\
$\downarrow$\\
Gazebo--PX4--ROS2闭环系统集成\\
$\downarrow$\\
统一场景、通信扰动与性能评价
\end{center}

\subsection{关键技术一：预测不确定性到运动基元风险的低复杂度映射}

重点解决动态障碍短时预测误差如何进入高频局部规划的问题。需要在预测均值、协方差或安全包络与候选运动基元之间建立计算开销可控的风险评价，使不确定性能够影响候选轨迹排序，同时与TTC和最小安全距离硬约束配合，避免增加风险表示后明显降低规划频率。

\subsection{关键技术二：通信可靠程度与邻机规划相关性的联合评价及约束调整}

重点解决通信状态和规划需求之间的对应关系。综合时延、丢包、AoI和链路状态描述共享信息可靠程度，再结合机间距离、相对运动、TTC和轨迹冲突趋势判断邻机规划相关性，据此筛选重点邻机，并动态调整轨迹可信时间窗、状态外推范围、机间安全裕度和冲突约束强度。

\subsection{关键技术三：多模块闭环集成与可重复实验接口}

重点解决Gazebo、PX4和ROS2之间多模块数据链路的一致性与可复现性问题，包括坐标系与状态定义统一、点云/障碍状态接口、规划指令仲裁、控制接口、多机命名空间、状态/轨迹消息共享、通信扰动注入及统一指标记录。通过模块化接口使算法、场景和通信参数能够独立替换，支持批量基线和消融实验。

\section{实验设计与评价方法}

\subsection{预测不确定性感知局部规划实验}

首先设置匀速、变速、转向以及带随机扰动的动态障碍运动模式，评价短时预测误差和不确定性表征效果，包括平均位移误差（ADE）、终点位移误差（FDE）、预测区间覆盖情况和单次预测计算耗时。随后将预测模块接入运动基元规划器，实验由未知静态场景逐步扩展到单动态障碍、多动态障碍及高密度动态场景。

设置三类核心方法进行比较：不使用未来预测的反应式局部规划、使用确定性未来预测的运动基元规划以及加入预测不确定性风险评价后的运动基元规划。评价任务成功率、碰撞率、最小安全距离、路径长度、轨迹平滑性、平均规划时间和规划频率，并通过障碍突然变速、转向或人为增加预测误差观察不同方法的性能退化情况。

进一步开展消融实验：去除不确定性风险项，仅保留确定性预测；保留风险项但去除TTC/安全距离硬过滤。实验不预设改进方法一定优于基线，而是重点判断其在预测误差增大时能否获得更稳定的安全表现，以及新增计算开销是否仍满足高频规划要求。

\subsection{通信约束多无人机协同实验}

构建多无人机动态环境，依次设置理想通信、固定时延、随机时延、不同丢包率、时延与丢包组合以及通信状态动态变化等场景。核心对比包括：（1）固定邻机集合与固定协同约束；（2）仅根据通信质量或信息新鲜度调整协同行为；（3）同时考虑通信状态与邻机规划相关性并调整协同约束的方法。

评价任务完成率、机间碰撞率、最小机间距离、任务完成时间、通信开销和规划计算时间。进一步通过去除邻机相关性筛选、保留邻机筛选但固定安全约束等消融方式，分析通信条件逐渐恶化时各部分对安全性、效率和通信/计算开销的实际贡献。

\subsection{闭环仿真系统验证实验}

系统验证按照“基础链路—单机动态—多机协同—通信受限综合场景”逐层进行。首先验证Gazebo传感器数据、ROS2规划节点、指令仲裁和PX4控制执行之间的闭环链路，包括起飞、悬停、航点跟踪、未知静态障碍避让及状态反馈；随后接入动态障碍状态估计、预测和风险感知规划，确认规划结果能够实时作用于飞行控制；再扩展到多无人机命名空间、状态/轨迹共享和协同避碰；最后注入通信时延、丢包和中断完成综合验证。

在同一套场景配置和指标统计接口下重复运行基线、改进方法和消融方案，记录碰撞、轨迹、控制状态、规划耗时、通信状态和异常情况。通过多组障碍密度、速度、机群规模和通信条件组合评价系统稳定性与实验可复现性，为前两个研究内容的结论提供统一实验支撑。

\section{可行性分析}

\subsection{理论基础可行}

动态障碍预测、不确定性风险建模、运动基元规划、多无人机分布式轨迹规划以及通信感知规划均已有较多研究基础\cite{gutow2020koopman,thomas2022sogm,wu2024sogm,robustMader}。本研究并非从零构建新的理论体系，而是在已有方法基础上针对高频局部规划中的预测不确定性利用方式，以及通信受限多机规划中的邻机筛选与协同约束调整开展研究，技术边界相对明确。

\subsection{算法实现可行}

第一项方法采用轻量局部规划框架，不依赖大规模端到端训练，短时预测可由状态估计和局部运动模型起步，不确定性可通过协方差传播、残差统计或安全包络表达。第二项方法主要位于多机协同管理和规划约束层，对PX4底层控制链路改动较小，可在已有单机规划基础上逐步扩展。

\subsection{实验平台可行}

Gazebo、PX4和ROS2能够分别提供场景与传感器仿真、无人机动力学和底层控制、上层感知与规划通信等能力。通过模块化接口可以逐步打通局部点云、规划指令、多机状态/轨迹共享和通信扰动链路，并通过参数化场景实现可重复实验。平台构建虽然工作量较大，但各模块可以分阶段验证，能够降低一次性集成风险。

\subsection{研究工作量与组织方式可控}

三个研究内容在论文表述上依次为“不确定性感知局部规划—通信受限协同规划—闭环仿真系统构建与综合验证”，其中前两个研究内容对应两项方法改进，第三个研究内容体现系统集成和实验工作量。实际实施时，基础平台链路先行并贯穿全过程：先支持第一项方法，再扩展支持第二项方法，最后整理形成统一的仿真验证系统。因此，研究内容顺序与工程实施顺序并不矛盾，且每个阶段均可形成独立的中间验证结果。
